\IfFileExists{ctexart.cls}{
  \documentclass[12pt,a4paper]{ctexart}
}{
  \documentclass[12pt,a4paper]{article}
}
\usepackage[left=2.5cm,right=2.5cm,top=2.8cm,bottom=2.8cm]{geometry}
\usepackage{setspace}
\usepackage{amsmath,amssymb}
\IfFileExists{booktabs.sty}{\usepackage{booktabs}}{
  % Fallbacks for minimal TeX installs
  \newcommand{\toprule}{\hline}
  \newcommand{\midrule}{\hline}
  \newcommand{\bottomrule}{\hline}
}
\IfFileExists{multirow.sty}{\usepackage{multirow}}{
  % Fallback: ignore multirow and render cell content normally
  \newcommand{\multirow}[3]{##3}
}
\usepackage{array}
\IfFileExists{hyperref.sty}{\usepackage[unicode]{hyperref}}{}
\usepackage{float}
\IfFileExists{longtable.sty}{\usepackage{longtable}}{}
\IfFileExists{caption.sty}{\usepackage{caption}}{}
\IfFileExists{enumitem.sty}{\usepackage{enumitem}}{}
\IfFileExists{natbib.sty}{\usepackage[numbers,sort&compress]{natbib}}{}
\usepackage{iftex}
\ifxetex
  \usepackage{fontspec}
  \setmainfont{Times New Roman}
  \IfFileExists{xeCJK.sty}{
    \usepackage{xeCJK}
    \IfFontExistsTF{PingFang SC}{
      \setCJKmainfont{PingFang SC}
    }{
      \IfFontExistsTF{Noto Serif CJK SC}{
        \setCJKmainfont{Noto Serif CJK SC}
      }{
        \IfFontExistsTF{Source Han Serif SC}{
          \setCJKmainfont{Source Han Serif SC}
        }{
          \IfFontExistsTF{Songti SC}{
            \setCJKmainfont{Songti SC}
          }{}
        }
      }
    }
  }{}
\fi

\hypersetup{
  colorlinks=true,
  linkcolor=blue,
  urlcolor=blue,
  citecolor=blue,
  unicode=true
}

\title{面向本体演化显著性评估的知识图谱存储自适应决策方法研究}
\author{作者：Yuliang}
\date{\today}

\begin{document}
\maketitle

\begin{abstract}
现有知识图谱与本体系统在演化场景下通常缺乏统一决策层：一方面，本体变化可能仅在编辑层发生而不影响查询收益；另一方面，盲目重构会引入额外重建成本、存储开销和稳定性风险。本文提出一种面向本体演化的存储自适应决策框架，将“变化显著性评估”和“代价驱动动作选择”统一到同一执行管线中。方法首先从相邻版本中提取变化事件，评估其结构影响、工作负载重叠和推理风险；随后生成候选动作集合，并通过效用函数在 no-op、refresh\_mv、repartition、reindex、cache\_plan 等动作中做选择。我们在 TPCH 与 PlantGraph 两类场景上完成五类实验：真实重放、决策质量、反事实验证、跨本体泛化和负载漂移。在当前实验设置下，所提方法在平均遗憾值（Average Regret）和累计净效用（CNU）上相对静态与单因素基线表现更优；在 exp1/exp2 中，AvgRegret 为 0.475、CNU 为 1090.2、Oracle 命中率为 0.75。本文同时给出可复现实验资产，支持脚本化复跑和论文产物自动汇总。
\end{abstract}

\textbf{关键词：}本体演化；显著性评估；存储自适应；动作选择；Regret 优化

\section{引言}
本体演化是知识图谱系统长期运行中的常态。已有研究分别在变化表示、版本管理、语义查询优化等方向形成了较成熟的方法体系\cite{flouris2008ontology,noy2004promptdiff,neumann2010rdf3x}，但在“是否触发重构、触发何种动作”这一在线决策问题上仍存在方法断层。实际系统中并非所有变化都值得触发代价较高的存储调整，且同一变化在不同工作负载下影响显著不同。

本文关注如下核心问题：\textbf{如何把本体变化从“编辑日志”转化为“可比较、可执行、可复现的存储决策信号”}。我们构建统一决策框架，将变化显著性评估作为触发条件，将代价函数优化作为动作选择依据，并通过统一 runner 保证各策略在同一输入输出口径下可公平对比。

\section{问题定义与研究问题}
\subsection{场景定义}
设本体版本序列为 $\{O_t\}_{t=1}^{T}$，相邻版本变化记为 $\Delta O_t = O_{t+1}-O_t$。系统在每个时刻观察到变化事件和工作负载状态，需在动作集合
\[
A=\{\text{no-op},\text{refresh\_mv},\text{repartition},\text{reindex},\text{cache\_plan}\}
\]
中选择动作 $\hat{a}_t$，以最大化长期效用。

\subsection{研究问题}
\begin{itemize}
  \item RQ1：如何定义本体变化显著性，使其反映对查询与推理的真实影响，而非仅反映 edit count？
  \item RQ2：如何将变化事件映射为可执行候选动作，并保证策略对比的一致性？
  \item RQ3：如何在查询收益、重建成本、存储开销和一致性风险间进行权衡？
  \item RQ4：在持续演化与负载漂移下，该方法是否优于静态和单因素基线？
\end{itemize}

\section{相关工作}
\subsection{本体演化与变化管理}
本体演化研究主要关注变化检测、版本维护与一致性保持。Flouris 等综述了本体变化管理的核心问题与框架\cite{flouris2008ontology}；Noy 与 Klein 提出的 PROMPTDIFF 在本体版本差分方面具有代表性\cite{noy2004promptdiff}。此类工作为“变化是什么”提供了方法基础，但通常不直接回答“变化是否值得触发存储动作”。

\subsection{RDF/知识图谱存储与查询优化}
在存储优化方面，RDF-3X 强调索引驱动的高效查询执行\cite{neumann2010rdf3x}，Hexastore 通过六向索引支持 RDF 三元组模式查询\cite{weiss2008hexastore}。这类方法主要针对查询执行效率优化，通常假设模式相对稳定，较少将本体演化信号直接纳入动作决策。

\subsection{工作负载基准与动态评测}
LUBM\cite{guo2005lubm} 与 WatDiv\cite{aluc2014watdiv} 是语义查询系统评估中的经典基准，分别覆盖规则场景与多样化查询结构。它们为负载对比提供了标准测试面，但未直接提供“变化事件到存储动作”的决策评估框架。本文在这些思路基础上，将演化事件、候选动作和 regret 指标整合到统一 runner 中进行比较。

\section{方法}
\subsection{总体框架}
方法流程为：\textbf{变化提取 $\rightarrow$ 显著性评估 $\rightarrow$ 候选动作生成 $\rightarrow$ 代价驱动选择}。其中显著性用于判断“该不该动”，代价模型用于回答“动什么”。

\subsection{显著性评估}
变化显著性定义为：
\[
S(\Delta O)=f(\text{structure},\text{workload},\text{reasoning},\text{constraints})
\]
分别表示结构影响、与热点工作负载重叠、推理扩张风险与约束敏感度。本文在工程实现中将其编码为事件特征（如 semantic\_significance、workload\_pressure 等）并输入统一策略接口。

\subsection{代价驱动动作选择}
动作效用定义如下：
\[
U(a)=\alpha \cdot \mathrm{QueryGain}(a)-\beta \cdot \mathrm{RebuildCost}(a)-\gamma \cdot \mathrm{StorageOverhead}(a)-\delta \cdot \mathrm{ConsistencyRisk}(a)
\]
最终选择
\[
\hat a=\arg\max_{a\in A(\Delta O)}U(a)
\]
Oracle 上界通过在候选动作集合内枚举真实效用获得，仅用于离线评估。

\section{实验设计}
\subsection{实验环境与可复现资产}
实验在项目统一脚本下执行，采用 Docker Compose 运行形态，TBOX 存储在 Neo4j，ABOX 存储在 Doris。可复现入口如下：
\begin{itemize}
  \item 预检查：\texttt{python experiments/preflight.py}
  \item 单项实验：\texttt{run\_exp1\_exp2.py}, \texttt{run\_exp3\_counterfactual.py}, \texttt{run\_exp4\_generalization.py}, \texttt{run\_exp5\_drift.py}
  \item 一键执行：\texttt{python experiments/run\_all.py}
\end{itemize}

\subsection{对比方法}
\begin{itemize}
  \item Static（从不触发）
  \item Always-Trigger（始终触发）
  \item Edit-Threshold（编辑阈值）
  \item Workload-Only（仅负载驱动）
  \item Significance-Only（仅显著性触发）
  \item Proposed（显著性+代价驱动）
  \item Oracle（离线上界）
\end{itemize}

\subsection{评价指标}
\paragraph{主指标}
\begin{align}
\mathrm{AvgRegret} &= \frac{1}{T}\sum_{t=1}^{T}\left(U(a_t^*)-U(\hat a_t)\right) \\
\mathrm{CNU} &= \sum_{t=1}^{T}U(\hat a_t) \\
\mathrm{UTR} &= \frac{\#\{t \mid \text{triggered and unnecessary}\}}{T} \\
\mathrm{HitRate} &= \frac{\#\{t \mid \hat a_t=a_t^*\}}{T}
\end{align}

\paragraph{辅指标}
P50/P95 延迟、重构耗时、存储开销、偏差率、一致性违例率。

\section{实验结果}
\subsection{实验一与实验二：真实重放与决策质量}
表~\ref{tab:exp12} 展示了连续版本重放场景中的主结果。在当前数据与参数设置下，Proposed 的 AvgRegret 最低（0.475），CNU 最高（1090.2），HitRate 为 0.75，相比主要基线具有更优表现。

\begin{table}[H]
\centering
\caption{Exp1/Exp2 结果汇总（12 steps）}
\label{tab:exp12}
\begin{tabular}{lcccc}
\toprule
Policy & AvgRegret $\downarrow$ & CNU $\uparrow$ & UTR $\downarrow$ & HitRate $\uparrow$ \\
\midrule
Static & 78.667 & 151.9 & 0.000 & 0.000 \\
Always-Trigger & 15.058 & 915.2 & 0.083 & 0.000 \\
Edit-Threshold & 22.783 & 822.5 & 0.000 & 0.000 \\
Workload-Only & 5.917 & 1024.9 & 0.000 & 0.583 \\
Significance-Only & 42.700 & 583.5 & 0.000 & 0.000 \\
\textbf{Proposed} & \textbf{0.475} & \textbf{1090.2} & 0.083 & \textbf{0.750} \\
\bottomrule
\end{tabular}
\end{table}

\subsection{实验三：反事实验证}
反事实结果显示，在“大变化低影响、小变化高影响、同 edit count 异影响”三类场景下，Proposed 与 Workload-Only 均可将 regret 控制在较低水平，其中 Workload-Only 的 AvgRegret 为 0.675，Proposed 为 0.850。该结果说明：仅依赖编辑次数不足以区分变化影响；同时在小样本条件下，策略间可能出现接近表现。

\subsection{实验四：跨本体泛化}
Leave-One-Ontology-Out 结果见表~\ref{tab:exp4}。无论 hold-out 为 PlantGraph 还是 TPCH，Proposed 在 AvgRegret 上均取得最低值（0.120 / 0.729），并保持较高 hit rate（0.8 / 0.714），显示出跨本体场景下的稳定性。

\begin{table}[H]
\centering
\caption{Exp4 跨本体泛化（LOOO）}
\label{tab:exp4}
\begin{tabular}{llccc}
\toprule
Held-out & Policy & AvgRegret $\downarrow$ & CNU $\uparrow$ & HitRate $\uparrow$ \\
\midrule
\multirow{2}{*}{PlantGraph} & Workload-Only & 9.920 & 447.1 & 0.600 \\
 & \textbf{Proposed} & \textbf{0.120} & \textbf{496.1} & \textbf{0.800} \\
\midrule
\multirow{2}{*}{TPCH} & Workload-Only & 3.057 & 577.8 & 0.571 \\
 & \textbf{Proposed} & \textbf{0.729} & \textbf{594.1} & \textbf{0.714} \\
\bottomrule
\end{tabular}
\end{table}

\subsection{实验五：负载漂移鲁棒性}
在 workload drift 场景下，Proposed 维持最低 AvgRegret（0.475）和最高 CNU（1090.2），Workload-Only 次之（AvgRegret 2.008，CNU 1071.8）。这表明联合信号在分布漂移条件下仍具备较好的收益稳定性。

\section{讨论}
\subsection{关键发现}
\begin{itemize}
  \item 仅依赖 edit count 或固定触发规则难以接近 Oracle。
  \item 将显著性与代价联合建模可显著降低长期 regret。
  \item 在跨域与漂移场景中，统一 runner 的口径优势能减少比较偏差。
\end{itemize}

\subsection{有效性威胁}
\begin{itemize}
  \item \textbf{内部有效性：}当前候选动作和效用映射仍是离散工程抽象，未来可引入真实线上执行反馈。
  \item \textbf{外部有效性：}当前主要覆盖 TPCH 和 PlantGraph，后续可扩展到 LUBM、WatDiv、行业本体。
  \item \textbf{构造有效性：}UTR 定义与业务阈值存在耦合，可在不同组织策略下做敏感性分析。
\end{itemize}

\section{结论与未来工作}
本文提出了面向本体演化场景的“显著性评估 + 代价驱动动作选择”统一方法，并在五类实验中进行了验证。基于当前实验结果，方法在长期演化过程中更接近 Oracle 决策并可获得更高净效用。未来工作将从三个方向推进：引入真实生产 trace 的在线学习、扩展动作空间（如分层物化与联邦缓存）、以及引入不确定性建模以支持风险约束优化。

\section*{附录：复现实验命令}
\begin{verbatim}
docker-compose up -d
bash scripts/init_doris.sh
bash scripts/seed_neo4j.sh
python experiments/run_all.py
\end{verbatim}

\begin{thebibliography}{99}
\bibitem{flouris2008ontology}
G. Flouris, D. Manakanatas, H. Kondylakis, D. Plexousakis, and G. Antoniou,
``Ontology change: classification and survey,''
\textit{The Knowledge Engineering Review}, vol. 23, no. 2, pp. 117--152, 2008.

\bibitem{noy2004promptdiff}
N. F. Noy and M. Klein,
``Ontology evolution: Not the same as schema evolution,''
\textit{Knowledge and Information Systems}, vol. 6, no. 4, pp. 428--440, 2004.

\bibitem{neumann2010rdf3x}
T. Neumann and G. Weikum,
``The RDF-3X engine for scalable management of RDF data,''
\textit{The VLDB Journal}, vol. 19, no. 1, pp. 91--113, 2010.

\bibitem{weiss2008hexastore}
C. Weiss, P. Karras, and A. Bernstein,
``Hexastore: sextuple indexing for semantic web data management,''
in \textit{Proceedings of the VLDB Endowment}, vol. 1, no. 1, pp. 1008--1019, 2008.

\bibitem{guo2005lubm}
Y. Guo, Z. Pan, and J. Heflin,
``LUBM: A benchmark for OWL knowledge base systems,''
\textit{Journal of Web Semantics}, vol. 3, no. 2--3, pp. 158--182, 2005.

\bibitem{aluc2014watdiv}
I. Aluc, O. Hartig, M. T. Ozsu, and K. Daudjee,
``Diversified stress testing of RDF data management systems,''
in \textit{The Semantic Web -- ISWC 2014}, pp. 197--212, 2014.
\end{thebibliography}

\end{document}

